\documentclass[a4paper,11pt,oneside]{book}
\usepackage{listings}
\usepackage[utf8]{inputenc}
\usepackage[spanish]{babel}
\usepackage{float}


\decimalpoint
\usepackage{dcolumn}
\newcolumntype{.}{D{.}{\esperiod}{-1}}
\makeatletter
\addto\shorthandsspanish{\let\esperiod\es@period@code}
\makeatother

\RequirePackage{verbatim}
\usepackage{fancyhdr}
\usepackage{graphicx}
\usepackage{afterpage}

\usepackage{longtable}

\usepackage[pdfborder={0 0 0}]{hyperref} % referencia


% ********************************************************************
% Re-usable information
% ********************************************************************
\newcommand{\myTitle}{Aplicación para la gestión de una tienda minorista\xspace}
\newcommand{\myDegree}{Grado en Ingeniería Informática\xspace}
\newcommand{\myName}{Julia María Cano Flores \xspace}
\newcommand{\myProf}{María Luisa Rodríguez Almendros \xspace}
\newcommand{\myFaculty}{Escuela Técnica Superior de Ingenierías Informática y de Telecomunicación\xspace}
\newcommand{\myFacultyShort}{E.T.S. de Ingenierías Informática y de Telecomunicación\xspace}
\newcommand{\myDepartment}{Departamento de Lenguajes y Sistemas Informáticos\xspace}
\newcommand{\myUni}{\protect{Universidad de Granada}\xspace}
\newcommand{\myLocation}{Granada\xspace}
\newcommand{\myTime}{\today\xspace}
\newcommand{\myVersion}{Version 0.1\xspace}

\hypersetup{
	pdfauthor = {\myName (email (en) ugr (punto) es)},
	pdftitle = {\myTitle},
	pdfsubject = {},
	pdfkeywords = {palabra_clave1, palabra_clave2, palabra_clave3, ...},
	pdfcreator = {LaTeX con el paquete ....},
	pdfproducer = {pdflatex}
}

\usepackage{url}
\usepackage{colortbl,longtable}
\usepackage[stable]{footmisc}
\usepackage{index}

\makeindex
%\usepackage[style=long, cols=2,border=plain,toc=true,number=none]{glossary}
%\makeglossary

\pagestyle{fancy}
\fancyhf{}
\fancyhead[LO]{\leftmark}
\fancyhead[RE]{\rightmark}
\fancyhead[RO,LE]{\textbf{\thepage}}
\renewcommand{\chaptermark}[1]{\markboth{\textbf{#1}}{}}
\renewcommand{\sectionmark}[1]{\markright{\textbf{\thesection. #1}}}


\usepackage{fancyhdr}
\pagestyle{fancy}
\fancyhf{} % Limpia el encabezado y pie de página de configuraciones previas
\fancyfoot[C]{\thepage} % Centra el número de página en el pie de página
\renewcommand{\headrulewidth}{0pt} % Elimina la línea del encabezado
\renewcommand{\footrulewidth}{0pt} % Elimina la línea del pie de página si existe

% Corrige cualquier espacio adicional alrededor del área del pie de página que podría desplazar el número de la página
\fancyfootoffset{0pt}

\usepackage{booktabs} % Para obtener líneas horizontales mejoradas
\usepackage[table]{xcolor}


\setlength{\headheight}{1.5\headheight}

% Configuraciones de colores y listados

\begin{document}
	% Portada
	\begin{titlepage}
 
 
\newlength{\centeroffset}
\setlength{\centeroffset}{-0.5\oddsidemargin}
\addtolength{\centeroffset}{0.5\evensidemargin}
\thispagestyle{empty}

\noindent\hspace*{\centeroffset}\begin{minipage}{\textwidth}

\centering
\includegraphics[width=0.9\textwidth]{imagenes/logo_ugr.jpg}\\[1.4cm]

\textsc{ \Large TRABAJO FIN DE GRADO\\[0.2cm]}
\textsc{ INGENIERÍA INFORMÁTICA}\\[1cm]
% Upper part of the page
% 
% Title
{\Huge\bfseries Aplicación para la gestión de una tienda minorista\\
}
\noindent\rule[-1ex]{\textwidth}{3pt}\\[3.5ex]

\end{minipage}

\vspace{2.5cm}
\noindent\hspace*{\centeroffset}\begin{minipage}{\textwidth}
\centering

\textbf{Autor}\\ {Julia María Cano Flores}\\[2.5ex]
\textbf{Directores}\\
{María Luisa Rodríguez Almendros}\\[2cm]
\includegraphics[width=0.3\textwidth]{imagenes/etsiit_logo.png}\\[0.1cm]
\textsc{Escuela Técnica Superior de Ingenierías Informática y de Telecomunicación}\\
\textsc{---}\\
Granada, febrero de 2024
\end{minipage}
%\addtolength{\textwidth}{\centeroffset}
%\vspace{\stretch{2}}
\end{titlepage}
\clearpage



	% Suprime la numeración de esta página
	\thispagestyle{empty}
	% Asegura que la numeración comience en la siguiente página
	\clearpage
	\setcounter{page}{2}
	
	% Prefacios y otros elementos preliminares
	\begin{titlepage}
 
 
\setlength{\centeroffset}{-0.5\oddsidemargin}
\addtolength{\centeroffset}{0.5\evensidemargin}
\thispagestyle{empty}

\noindent\hspace*{\centeroffset}\begin{minipage}{\textwidth}

\centering
%\includegraphics[width=0.9\textwidth]{imagenes/logo_ugr.jpg}\\[1.4cm]

%\textsc{ \Large PROYECTO FIN DE CARRERA\\[0.2cm]}
%\textsc{ INGENIERÍA EN INFORMÁTICA}\\[1cm]
% Upper part of the page
% 

 \vspace{3.3cm}

%si el proyecto tiene logo poner aquí
\includegraphics{imagenes/logo.png} 
 \vspace{0.5cm}

% Title

{\Huge\bfseries Aplicación para la gestión de una tienda minorista\\
}
\noindent\rule[-1ex]{\textwidth}{3pt}\\[3.5ex]

\end{minipage}

\vspace{2.5cm}
\noindent\hspace*{\centeroffset}\begin{minipage}{\textwidth}
\centering

\textbf{Autor}\\ {Julia María Cano Flores}\\[2.5ex]
\textbf{Directores}\\
{María Luisa Rodríguez Almendros}\\[2cm]
\includegraphics[width=0.15\textwidth]{imagenes/tstc.png}\\[0.1cm]
\textsc{Departamento de Lenguajes y Sistemas Informáticos}\\
\textsc{---}\\
Granada, febrero de 2024
\end{minipage}
\addtolength{\textwidth}{\centeroffset}
\vspace{\stretch{2}}

 
\end{titlepage}





\cleardoublepage
%\thispagestyle{empty}

\begin{center}
{\large\bfseries Aplicación para la gestión de una tienda minorista}\\
\end{center}
\begin{center}
Julia María Cano Flores\\
\end{center}

\vspace{0.7cm}
\noindent{\textbf{Palabras clave}: minorista, base de datos, software, requisitos}\\

\vspace{0.7cm}
\noindent{\textbf{Resumen}}\\

El desarrollo de este Trabajo Fin de Grado (TFG) tiene como finalidad principal simplificar y mejorar la eficiencia de la gestión de negocios minoristas. Está diseñado para optimizar el proceso de seguimiento de ventas mediante una base de datos que registre todas las transacciones realizadas. El sistema no solo registra compras y devoluciones, sino que también incorpora una característica distintiva: la gestión de préstamos. Esta funcionalidad es particularmente útil en localidades pequeñas, donde las relaciones de confianza entre comerciantes y clientes permiten prácticas como llevarse productos sin un desembolso inicial. Los clientes pueden probar los artículos antes de comprometerse a la compra o optar por la devolución sin costos adicionales.\\


El software proporcionará una interfaz de usuario intuitiva compuesta de múltiples pantallas, entre las que podemos destacar un inventario actualizado en tiempo real, la visualización de los artículos en venta o una lista de clientes habituales del negocio. Además, para impulsar la toma de buenas decisiones estratégicas, el programa generará gráficos analíticos que reflejarán el progreso económico del negocio. \\

Este proyecto surge de la necesidad evidente de una solución tecnológica adaptada a los requisitos específicos de los negocios minoristas. Hasta el momento, son negocios que no se han tenido en cuenta debido a su bajo impacto y su pequeño tamaño. Sin embargo, debemos de mirar por llevar la tecnología a todos los ámbitos. Esta aplicación busca llenar ese vacío y convertirse en una herramienta útil para comerciantes minoristas.  


\cleardoublepage


%\thispagestyle{empty}


\begin{center}
{\large\bfseries Application for the management of a retail store}\\
\end{center}
\begin{center}
Julia María Cano Flores\\
\end{center}

\vspace{0.7cm}
\noindent{\textbf{Keywords}: retailer, database, software, requirements}\\

\vspace{0.7cm}
\noindent{\textbf{Abstract}}\\


The main purpose of the development of this Final Degree Project (FDP) is to simplify and improve the efficiency of retail business management. It is designed to optimize the sales tracking process through a database that records all transactions made. The system not only records purchases and returns but also incorporates a distinctive feature: loan management. This functionality is particularly useful in small localities, where trust relationships between merchants and customers allow practices such as taking products without an initial outlay. Customers can test items before committing to purchase or opt for a return without additional costs.\\

The software will provide an intuitive user interface composed of multiple screens, among which we can highlight an inventory updated in real-time, the display of items for sale, or a list of regular customers of the business. In addition, to drive good strategic decision-making, the program will generate analytical charts that will reflect the economic progress of the business.\\

This project arises from the clear need for a technological solution adapted to the specific requirements of retail businesses. So far, they are businesses that have not been considered due to their low impact and small size. However, we must look to bring technology to all areas. This application seeks to fill that void and become a useful tool for retail merchants.



\chapter*{}

%\thispagestyle{empty}


\noindent\rule[-1ex]{\textwidth}{2pt}\\[4.5ex]

Yo, \textbf{Julia María Cano Flores}, alumna de la titulación INGENIERÍA INFORMÁTICA de la \textbf{Escuela Técnica Superior
de Ingenierías Informática y de Telecomunicación de la Universidad de Granada}, con DNI 77649643x, autorizo la
ubicación de la siguiente copia de mi Trabajo Fin de Grado en la biblioteca del centro para que pueda ser
consultada por las personas que lo deseen.

\vspace{6cm}

\noindent Fdo: Julia María Cano Flores

\vspace{2cm}

\begin{flushright}
Granada a 2 de febrero de 2024 .
\end{flushright}


\chapter*{}
%\thispagestyle{empty}

\noindent\rule[-1ex]{\textwidth}{2pt}\\[4.5ex]

D. \textbf{María Luisa Rodríguez Almendros}, Profesora del Área de Ingeniería del Software del Departamento de Lenguajes y Sistemas Informáticos de la Universidad de Granada.




\vspace{0.5cm}

\textbf{Informan:}

\vspace{0.5cm}

Que el presente trabajo, titulado \textit{\textbf{Aplicación para la gestión de una tienda minorista}},
ha sido realizado bajo su supervisión por \textbf{Julia María Cano Flores}, y autorizamos la defensa de dicho trabajo ante el tribunal
que corresponda.

\vspace{0.5cm}

Y para que conste, expiden y firman el presente informe en Granada a 2 de febrero de 2024 .

\vspace{1cm}

\textbf{Los directores:}

\vspace{5cm}

\noindent \textbf{María Luisa Rodríguez Almendros}

\chapter*{Agradecimientos}
\thispagestyle{empty}

       \vspace{1cm}


Poner aquí agradecimientos...




	
	\tableofcontents
	\listoffigures
	\listoftables
	
%	\mainmatter
	\setlength{\parskip}{5pt}
	
	% Contenido principal
	\chapter{Introducción}
\label{chap:introduction}

\section{Motivación}
\label{sec:motivation}

Se entiende como negocio minorista toda empresa de comercio que adquiere mercancías por cuenta propia, y las revende directamente al consumidor final. En negocios minoristas pequeños podemos encontrar un trato más cercano con los clientes y una forma algo distinta de gestionar un comercio. \\

Al ser nacida y criada en un pueblo, la mayoría de negocios con los que he crecido eran comercios minoristas pequeños. Además, mi padre es propietario de una tienda y he podido conocer de primera mano cuál es el funcionamiento real de este tipo de negocios. Tras una vida tratando con ellos y una entrevista realizada a mi padre, he encontrado la necesidad de informatizar este trabajo con el objetivo de conseguir optimizar el rendimiento y llevar un seguimiento. En la actualidad, muchos comerciantes minoristas no disponen de grandes tecnologías y recurren a métodos tradicionales basados en papel para el seguimiento de su negocio. Esto tiene un riesgo alto de pérdida de datos o cometer errores. 
Además, el proceso de recuperación y análisis de datos se convierten en tareas muy tediosas, ya que la falta de sistemas automatizados impide un filtrado eficiente de información. Todo esto ralentiza el trabajo de forma significativa, haciendo que la productividad del comerciante sea menor y, por tanto, la evolución del negocio no sea óptima. \\

En el mercado actual no encontramos aplicaciones que pongan solución al problema de los negocios minoristas pequeños. Las soluciones software existentes tienden a enfocarse en entidades de mayor escala. Como hemos visto, los negocios minoristas tienen singularidades que conllevan establecer requisitos específicos y tratar su gestión de forma independiente a la gestión de un gran negocio. Por ello, es crucial el desarrollo de un software que satisfaga las necesidades de los pequeños comerciantes y les permita conducir su negocio a una evolución óptima. 


\section{Objetivos del proyecto}
\label{sec:project_objectives}

El objetivo general de este proyecto es proporcionar una solución informática analizando los procedimientos típicos de un comercio minorista pequeño que actualmente se realizan de forma manual y desarrollando una aplicación que los automatice. Para conseguir dicho objetivo, podemos diferenciar una serie de objetivos específicos: 


% Tabla para el Objetivo 1
\definecolor{grayshade}{gray}{0.9}



\begin{itemize}
	\item \textbf{Objetivo 1: Análisis de los comercios minoristas}
	\begin{itemize}
		\item \textbf{Descripción:} Se debe de realizar una búsqueda de información sobre los negocios minoristas con el objetivo de entender mejor cuáles son sus características.
	\end{itemize}
	
	\item \textbf{Objetivo 2: Estudio de aplicaciones similares}
	\begin{itemize}
		\item \textbf{Descripción:} Análisis del mercado actual de aplicaciones similares para identificar funcionalidades innovadoras que implementar en nuestra aplicación.
	\end{itemize}
	
	\item \textbf{Objetivo 3: Análisis de usabilidad y accesibilidad}
	\begin{itemize}
		\item \textbf{Descripción:} Estudiar soluciones accesibles para conseguir una aplicación fácil de usar para la mayoría de usuarios.
	\end{itemize}
	
	\item \textbf{Objetivo 4: Análisis de las tecnologías a utilizar}
	\begin{itemize}
		\item \textbf{Descripción:} Investigar sobre cuáles son las mejores tecnologías para llevar a cabo el proyecto.
	\end{itemize}
	
	\item \textbf{Objetivo 5: Desarrollo e implementación de la aplicación}
	\begin{itemize}
		\item \textbf{Descripción:} Analizar, desarrollar y probar la aplicación con el objetivo de conseguir los mejores resultados.
	\end{itemize}
	
	\item \textbf{Objetivo 6: Validación de la aplicación por usuarios reales}
	\begin{itemize}
		\item \textbf{Descripción:} Poner a prueba la aplicación en un entorno real para ver si cumple las expectativas del usuario.
	\end{itemize}
\end{itemize}

\section{Estructura del documento}
\label{sec:document_structure}
Esta sección se ofrece una visión de la estructura y organización del documento. Se compone de las siguientes partes: 

\begin{itemize}
	\item \textbf{Resumen:} Un resumen donde se exponen las principales funcionalidades del proyecto y sus objetivos. 
	
	\item \textbf{Introducción:} En este apartado se exponen los principales objetivos que se pretenden cumplir durante el desarrollo del proyecto. Además, se explica la motivación por la que se empezó a idear dicho proyecto a modo de introducción. 

	
	\item \textbf{Planificación:} Especificación de la planificación del proyecto, donde se muestran las distintas fases y la estimación presupuestaria del mismo. 
	
	\item \textbf{Estado del arte:} Investigación sobre los principales temas que se tratan en el proyecto como los negocios minoristas, los negocios mixtos y los comerciantes minoristas. Además, se tratan temas legales como la protección de datos. Con esta investigación, entramos en contexto para posteriormente entender los requisitos del proyecto. 
	
	\item \textbf{Análisis de tecnologías a utilizar:} Enumeración de tecnologías utilizadas para el desarrollo del proyecto. Para cada tecnología, se explica porqué se ha utilizado y se citan sus principales características. 
	
	\item \textbf{Análisis de la aplicación:} Obtención y explicación de los requisitos del sistema. Posteriormente se amplia esta información desarrollando los casos de uso, el modelo conceptual y los diagramas de secuencia del sistema. 
	
	\item \textbf{Diseño:} Descripción de la arquitectura del sistema, el diseño de las clases necesarias para el programa y la interfaz de usuario.
	
	\item \textbf{Implementación de la aplicación:} La implementación de la aplicación se desarrollará en con una metodología ágil, por lo que este capítulo se dividirá en sus iteraciones. 
	
	\item \textbf{Bibliografía:} Recopilación de todas las fuentes de información utilizadas durante la realización del proyecto.
\end{itemize}
\newpage






	
\end{document}



%
%\input{capitulos/02_EspecificacionRequisitos}
%
%\input{capitulos/03_Planificacion}
%
%\input{capitulos/04_Analisis}
%
%\input{capitulos/05_Diseno}
%
%\input{capitulos/06_Implementacion}
%
%\input{capitulos/07_Pruebas}
%
%\input{capitulos/08_Conclusiones}
%
%%\chapter{Conclusiones y Trabajos Futuros}
%
%
%%\nocite{*}
%\bibliography{bibliografia/bibliografia}\addcontentsline{toc}{chapter}{Bibliografía}
%\bibliographystyle{miunsrturl}
%
%\appendix
%\input{apendices/manual_usuario/manual_usuario}
%%\input{apendices/paper/paper}
%\input{glosario/entradas_glosario}
% \addcontentsline{toc}{chapter}{Glosario}
% \printglossary
%\chapter*{}
%\thispagestyle{empty}


